\section{Stage 2: Multi-Source Candidate Ranking}
Once candidate data and a job description have been extracted and combined in a configuration, we can then proceed to
the execution stage. In this stage, we compute signals for each metric. This can be done in two different ways
depending on whether the metric is baseline or extensible. Both metrics may take into account multiple data sources,
computing a value (signal) per data source where relevant. When a metric has more than one signal, we perform Bayesian 
fusion to compute a single score. Once each metric has been reduced to a single value, we then undergo a 
weighted average of all the metrics to compute the final (match) score. Figure \ref{fig:schema_overview} 
shows how we can use a candidate's profile alongside a job description to produce a final score: candidates from multiple 
sources are grouped into a batch, paired with an extracted job description to form a configuration, and each configuration 
run stores an immutable past result.


\begin{figure}[htbp]
    \centering
    \begin{tikzpicture}[
        node distance=1.5cm and 2.5cm,
        arch_table/.style={rectangle, draw=CoolGrey, fill=white, text width=3.2cm, align=left, font=\footnotesize\sffamily, rounded corners=1pt, drop shadow, minimum height=1.6cm},
        arch_source/.style={ellipse, draw=ImperialBlue, fill=LightBlue!40, font=\footnotesize\bfseries\sffamily, minimum width=2.0cm, align=center},
        arch_group/.style={draw=ImperialBlue, fill=LightBlue!10, dashed, rounded corners, inner sep=10pt},
        pk/.style={font=\footnotesize\bfseries\sffamily, ImperialBlue},
        fk/.style={font=\footnotesize\itshape\color{ProcessBlue}},
        arrow/.style={-{Stealth[scale=1.0]}, thick, draw=ImperialBlue}
    ]
        % Candidate Schema Part: Sources in a column on the left
        \node [arch_table] (t_gh) {
            \textbf{github\_profiles}\\
            {\color{ImperialBlue}\bfseries id}\\
            username, total\_commits\\
            total\_stars, languages
        };

        \node [arch_table, below=0.6cm of t_gh] (t_li) {
            \textbf{linkedin\_profiles}\\
            {\color{ImperialBlue}\bfseries id}\\
            full\_name, headline\\
            location, about
        };

        \node [arch_table, below=0.6cm of t_li] (t_edu) {
            \textbf{candidate\_education}\\
            {\color{ImperialBlue}\bfseries id}\\
            {\itshape\color{ProcessBlue} candidate\_id}\\
            school\_name, degree, grade
        };

        % Central Identity Table (Closer to sources)
        \node [arch_table, right=1.6cm of t_li] (t_cand) {
            \textbf{candidate\_data}\\
            {\color{ImperialBlue}\bfseries id}\\
            name, email, skills\\
            {\itshape\color{ProcessBlue} github\_profile\_id}\\
            {\itshape\color{ProcessBlue} linkedin\_profile\_id}
        };

        % Configuration Schema Part: Anchored further right to avoid collision
        \node [arch_table, right=6cm of t_gh] (t_jd) {
            \textbf{job\_requirements}\\
            {\color{ImperialBlue}\bfseries id}\\
            title, description\\
            metrics (JSONB)
        };

        \node [arch_table, below=0.6cm of t_jd] (t_conf) {
            \textbf{matching\_configs}\\
            {\color{ImperialBlue}\bfseries id}\\
            {\itshape\color{ProcessBlue} job\_id}, {\itshape\color{ProcessBlue} batch\_id}\\
            weights, active\_metrics
        };

        % Results
        \node [arch_table, below=0.6cm of t_conf] (t_res) {
            \textbf{past\_results}\\
            {\color{ImperialBlue}\bfseries id}\\
            {\itshape\color{ProcessBlue} config\_id}\\
            results\_payload (JSONB)
        };

        % Groupings
        \begin{scope}[on background layer]
            \node [arch_group, fit=(t_gh) (t_li) (t_edu) (t_cand), label={[font=\footnotesize\bfseries\sffamily, ImperialBlue]above:CANDIDATE PERSISTENCE}] (db_cand) {};
            \node [arch_group, fit=(t_jd) (t_conf) (t_res), label={[font=\footnotesize\bfseries\sffamily, ImperialBlue]above:CONFIG \& RESULTS}] (db_conf) {};
        \end{scope}

        % Flows
        \draw [arrow] (t_jd) -- (t_conf);
        \draw [arrow] (t_conf) -- (t_res);
        
        % Internal Relationships
        \draw [arrow] (t_gh.east) -- ++(0.4,0) |- (t_cand.west);
        \draw [arrow] (t_li.east) -- (t_cand.west);
        \draw [arrow] (t_edu.east) -- ++(0.4,0) |- (t_cand.west);

        % Flow to Matching (Clear crossing between groups)
        \draw [arrow] (t_cand.east) -- (t_conf.west);

    \end{tikzpicture}
    \caption{Relational Schema: Integration of candidate signals, recruiter criteria, 
    and auditable matching results.}
    \label{fig:schema_overview}
\end{figure}



\subsection{Semantic Skill Matching}
\label{subsec:semantic-skill-matching}

One of the biggest issues with strict keyword-matching screening systems lies in their rigidity.
Suppose a candidate applied to a role where ``SQL'' is required, however, the candidate only uses
``PostgreSQL'' or ``Postgres'' in their CV. This candidate would likely achieve a zero for the ``SQL'' metric.
Semantic skill matching (or bridging) addresses this issue. When there are no exact matches,
we perform a cosine similarity operation, using the \textit{SentenceTransformers} library with the 
\textit{all-MiniLM-L6-v2} model. We selected this model due to its efficiency and accuracy regarding
technical terminology. 

For this model, we encode the skills into a 384-dimensional vector space, and then compute cosine
similarity (Equation~\ref{eq:cosine_similarity}).

To be considered as a semantic bridge, there must be a cosine similarity score of $\geq 0.50$ between
the extensible metric and the skill mentioned in the candidate CV. This threshold was determined 
through iterative testing. The threshold is permissive enough to match terms such as "SQL" and "PostgreSQL",
while also remaining strict enough to reject other, unrelated terms such as ``Python'' and ``JavaScript''.
Whenever there is a semantic bridge, we also pass this information to the frontend to let the recruiter know
of the extensible metric, as well as the skill that was bridged against it, to ensure full transparency.
This way, in the event of a false-positive, the recruiter can easily identify and correct it.
For extensible CV metrics, a direct string match (e.g.\ ``SQL'') yields a lexical signal, otherwise cosine similarity
against extracted skills is used (e.g.\ bridging ``PostgreSQL'' to ``SQL''), as summarised in Figure~\ref{fig:semantic_skill_scoring}.


\begin{figure}[htbp]
    \centering
    \resizebox{0.98\textwidth}{!}{%
    \begin{tikzpicture}[
        arch_block/.style={
            rectangle, draw=ImperialBlue, fill=white, text width=1.65cm,
            align=center, minimum height=0.62cm, rounded corners=2pt,
            font=\scriptsize\sffamily, line width=0.7pt
        },
        arch_source/.style={
            rectangle, draw=CoolGrey, fill=LightBlue!20, text width=1.5cm,
            align=center, minimum height=0.55cm, font=\scriptsize\sffamily,
            line width=0.7pt
        },
        arch_math/.style={
            diamond, draw=ImperialBlue, fill=LightBlue!10, text width=0.95cm,
            align=center, inner sep=0pt, font=\tiny\sffamily, line width=0.7pt,
            aspect=1.45
        },
        arch_decision/.style={
            diamond, draw=ImperialBlue, fill=white, text width=1.05cm,
            align=center, inner sep=0pt, font=\tiny\sffamily, line width=0.7pt,
            aspect=1.45
        },
        arr/.style={
            -{Stealth[scale=1.0]}, thick, draw=ImperialBlue
        },
        busline/.style={
            thick, draw=ImperialBlue
        },
        edgelabel/.style={font=\tiny\itshape, inner sep=2pt, fill=white}
    ]
        % Row centres at y = ±1.05; columns spaced on a fixed grid
        \node [arch_source] (metric) at (0, 1.05) {\textbf{JD Metric}\\{\tiny e.g.\ PostgreSQL}};
        \node [arch_source] (skills) at (0, -1.05) {\textbf{CV Skills}\\{\tiny canonical}};

        \node [arch_decision] (exact) at (3.5, 0) {Exact\\match?};

        \node [arch_block] (lex) at (6.2, 1.05) {\textbf{Lexical}\\CV signal};
        \node [arch_block] (sbert) at (6.2, -1.05) {\textbf{SBERT}\\{\tiny all-MiniLM-L6-v2}};
        \node [arch_math] (cosine) at (8.5, -1.05) {Cosine\\{\tiny $\geq 0.50$}};

        \node [arch_block] (out) at (13.2, 0) {\textbf{CV metric}\\signal};

        % Shared vertical buses (schema-style joints)
        \coordinate (inBus) at (1.65, 0);
        \coordinate (branchBus) at (4.55, 0);
        \coordinate (mergeBus) at (11.0, 0);

        % Inputs tap a vertical bus (horizontal arrows only; no opposing heads)
        \draw [arr] (metric.east) -- (inBus |- metric.east);
        \draw [arr] (skills.east) -- (inBus |- skills.east);
        \draw [busline] (inBus |- metric.east) -- (inBus |- skills.east);
        \draw [arr] (inBus) -- (exact.west);

        % yes / no on horizontal legs (symmetric from branch bus)
        \draw [arr] (exact.east) -- (branchBus);
        \draw [arr] (branchBus) -- (branchBus |- lex.east)
            node[pos=0.55, above, edgelabel] {yes} -- (lex.west);
        \draw [arr] (branchBus) -- (branchBus |- sbert.east)
            node[pos=0.55, below, edgelabel] {no} -- (sbert.west);

        % Semantic path (straight, same row as SBERT)
        \draw [arr] (sbert.east) -- (cosine.west);

        % Branches tap merge bus (horizontal in; one arrow out)
        \draw [arr] (lex.east) -- (mergeBus |- lex.east);
        \draw [arr] (cosine.east) -- (mergeBus |- cosine.east);
        \draw [busline] (mergeBus |- lex.east) -- (mergeBus |- cosine.east);
        \draw [arr] (mergeBus) -- (out.west);

    \end{tikzpicture}%
    }
    \caption[Stage~2: semantic skill bridging]{Semantic skill bridging during the scoring of
    extensible metrics for the CV. If we see a direct match on an extensible metric (E.g. "SQL"),
    we have a lexical CV signal, similar to keyword matching. If there is not a direct match (E.g. "PostgreSQL"),
    we perform cosine similarity between the skills we have extracted, and the metric in the configuration.
    If there is a match ($\geq 0.50$), then there is a valid semantic bridge, and we then update the CV metric signal.}
    \label{fig:semantic_skill_scoring}
\end{figure}


\subsection{Integrity Layer}
A common flaw in legacy ATS platforms lies in their vulnerability to adversarial attacks, such as
keyword-stuffing or identity squatting. To combat these, MERIT executes an Integrity Audit as a strict
pre-processing step immediately prior to signal evaluation. While the detailed heuristics for
these safeguards are discussed in Section~\ref{sec:bias-mitigation-safeguards}, 
it is important to note that the application of these penalties depends on the nature of the attack. 
For signal-specific attacks like keyword-stuffing, the audits are calculated first, allowing the resulting penalties to be
injected directly into the raw evidence signals before they enter the Bayesian fusion pipeline. This ensures that
poisoned data carries less significance in the final probabilistic model, meaning a candidate may
still rank highly if they have strong, authentic evidence from other sources. Conversely, for severe global threats 
like ``Identity Squatting'' (where a candidate passes off another developer's profile as their own), the penalty circumvents 
the Bayesian fusion process entirely and applies a severe, flat deduction to the candidate's final aggregate score. 
Both of these mechanisms are covered in more detail in Section~\ref{sec:bias-mitigation-safeguards}.

\subsection{Skill Decay Modelling}
\label{subsec:skill-decay-modelling}
Skill decay refers to a model that predicts the loss of a skill over time due to the lack of use.
Applying this to technical recruitment, it is intuitive that the recency of a candidate's skill matters. A candidate may have been considered
an expert in a certain programming language ten years ago, but if they had not used it since then, would their level be considered the same?
They may not carry the same weight as a candidate who has used it often within the past few months, and prior research shows that this skill
loss grows with longer periods of non-use \cite{arthur1998}. Industry research points to a similar finding, where workplace skills (especially technical skills)
become outdated far more quickly than they once did \cite{ibm2019skillsgap,wef2020futurejobs}. Traditional ATS platforms fail to account for this since
there is often no significant evidence to uncover when the candidate last used the skill in question.

MERIT attempts to overcome this issue through the addition of an exponential skill decay model for extensible metrics (e.g.\ ``Python'').
We apply this to the GitHub signal, as it is the most reliable signal for demonstrating evidence of a candidate's skill.
For other signals, it is usually harder to find a point to start the decay from, but for GitHub, we can determine it from the candidate's historical commits.
In our backend, we extract the repository's language distribution and compute a volume-weighted average across its active years.
This calculates the candidate's effective ``center of mass'' for that skill, which is far more robust against candidates pushing tiny, 
superficial commits to legacy repositories to simulate recent activity. Subtracting this weighted year from the current
date yields the $\Delta t$ value in the decay formula:

\begin{equation} \label{eq:exponential_decay}
  \mathcal{S}_{\text{decay}} = \mathcal{S}_{\text{base}} \cdot e^{-\lambda \Delta t}
\end{equation}

Where:
\begin{itemize}
  \item $\mathcal{S}_{\text{base}}$ is the initial magnitude of the extracted skill signal.
  \item $\Delta t$ is the time elapsed (in years) since the end date of the relevant activity.
  \item $\lambda$ is the decay constant, determining the rate of skill degradation.
\end{itemize}

The system is configured with a default exponential decay constant of $\lambda = 0.25$,
which corresponds to a half-life of $t_{1/2} = \ln(2)/\lambda \approx 2.8$ years.
This value is not taken directly from a single empirical study. Instead, it is a practical default chosen to align with
industry estimates that technical skills can lose much of their relevance within roughly 2.5--5 years \cite{ibm2019skillsgap}.
This default value ensures that recent, relevant experience naturally outranks historical,
potentially obsolete capabilities.

\begin{figure}[htbp]
    \centering
    \begin{tikzpicture}[
            font=\sffamily,
            xscale=1.1, yscale=4.0
        ]
        % Origin
        \node[below left, font=\scriptsize] at (0,0) {0};

        % Grid lines
        \foreach \x in {1,2,3,4,5,6,7,8,9,10} {
                \draw[gray!30, dashed] (\x,0) -- (\x,1);
                \node[below, font=\scriptsize] at (\x,0) {\x};
            }
        \foreach \y in {0.2, 0.4, 0.6, 0.8, 1.0} {
                \draw[gray!30, dashed] (0,\y) -- (10,\y);
                \node[left, font=\scriptsize] at (0,\y) {\y};
            }

        % Axes
        \draw[thick, -{Stealth[scale=1.0]}] (0,0) -- (10.5,0) node[right, font=\footnotesize\bfseries] {Years Inactive ($\Delta t$)};
        \draw[thick, -{Stealth[scale=1.0]}] (0,0) -- (0,1.15) node[above, font=\footnotesize\bfseries] {Signal Strength ($\mathcal{S}_{\text{decay}}$)};

        % Decay Curve: y = e^(-0.25 * x)
        \draw[ultra thick, ImperialBlue, domain=0:10, samples=100]
        plot (\x, {exp(-0.25 * \x)});

        % Half-life annotation
        \draw[dashed, thick, CoolGrey] (2.77,0) -- (2.77,0.5) -- (0,0.5);
        \node[circle, fill=CoolGrey, inner sep=1.5pt] at (2.77,0.5) {};
        \node[above right, font=\footnotesize\bfseries] at (2.77,0.5) {Half-Life Point (2.8 Years)};

        % Start point
        \node[circle, fill=ImperialBlue, inner sep=1.5pt] at (0,1.0) {};

    \end{tikzpicture}
    \caption[Temporal skill decay ($\lambda = 0.25$)]{Temporal skill decay with $\lambda = 0.25$ (half-life $\approx 2.8$ years).}
    \label{fig:skill_decay_graph}
\end{figure}




\subsection{Separation of Evidence Types}
The reasoning for MERIT's multi-source design is to be able to consider different sources of candidate data,
those that are self-reported by the candidate, and those that the candidate has actually been able to demonstrate
proficiency in. MERIT splits these into two different forms of evidence, \textbf{Claimed Evidence} and 
\textbf{Demonstrated Evidence}.

\textbf{Claimed evidence} refers to a self-declared proficiency in a skill, programming language or framework.
This form of evidence is commonly found in CVs and LinkedIn profiles due to the user being responsible
for their own claims. This is the only form of evidence usually considered by ATS scoring engines, making
these systems biased towards candidates who are able to claim their skills in a more favourable light. Self-embellishment
is often never punished, as it is difficult to verify, but this is where the beauty of MERIT comes into play,
with its second form of evidence.

\textbf{Demonstrated Evidence} refers to evidence derived from auditable, historical actions. In our case, this is 
specifically GitHub repository metadata. For example, a candidate is likely to claim Python proficiency on their CV for 
a job application, however, this claim can be challenged by the presence of this skill on their GitHub repositories.
Demonstrated evidence is significantly harder to game as it requires a candidate to actually use the skill as part of 
their own projects or daily work.

This separation is not just conceptual, this is a formal architectural separation. By splitting data sources into claimed
and demonstrated sources, MERIT can assign different levels of trust to these sources, which will play a key role in the 
Bayesian Fusion algorithm discussed next.

\subsection{Bayesian Fusion}
\label{subsec:bayesian-fusion}
The most significant innovation in MERIT's scoring engine is the adoption of \textbf{Bayesian Evidence Fusion}.
Initial iterations of MERIT attempted to combine skill signals using static weighted averages. However, this approach often dilutes
the importance of demonstrated evidence (GitHub) with claimed evidence (CV and LinkedIn). This approach also failed to account for the
confidence of a measurement. For example, a single line in a CV stating "Python" should not carry the same weight as a GitHub profile
that demonstrates extensive Python experience.

Another key issue is the lack of flexibility it provides, as we would now need separate definitions to account for candidates who may not
have all three data sources. Several considered approaches failed to provide a fair solution whilst adhering to the static weighted average
approach. Alongside this, there was also the issue of extensibility, as we need to provide $2^N - 1$ distinct definitions for $N$ data sources if
we wish to avoid penalising candidates for missing data sources.

The presence of these dilemmas led to the integration of Bayesian Evidence Fusion. MERIT now treats each
data source as independent signals updating a Beta Distribution. We parameterise this with the use of a sparse prior
($\alpha_{prior} = 0.1, \beta_{prior} = 0.1$) to ensure high sensitivity to initial evidence signals while maintaining a baseline for uncertainty.

The contribution of each source $s$ to the final score of that metric is given by its raw signal $\sigma_s \in [0, 1]$ and its reliability weight $w_s$.
Evidence-based data sources are often assigned a higher weight due to MERIT's emphasis on evidence-based screening. To encourage candidates to
use evidence-based sources while avoiding harsh penalisation, we enforce a ``Signal Cap'' for sources that provide claimed evidence.
We then map these signals to the evidence parameters as follows:
\begin{equation}
  \alpha_s = w_s \sigma_s + \epsilon_s, \quad \beta_s = w_s (1 - \sigma_s) + \epsilon_s
\end{equation}
where $\epsilon_s = 0.01(1 - w_s)$ represents a symmetric uncertainty term. The reason for the inclusion of $\epsilon_s$ is to inject noise for
low-confidence sources to produce wider, more uncertain distributions. This stops the distribution from becoming narrow unless this candidate has the
evidence to support their claims (since high-confidence sources provide more concentrated signals). Crucially, this noise injection adheres to Cromwell's Rule, mathematically ensuring the engine never reaches absolute certainty ($1.0$) since digital footprints cannot provide omniscient proof of a candidate's abilities.

We select the Beta distribution since it is the conjugate prior for Bernoulli-distributed evidence \cite{gelman2013bayesian}. This means that the resulting posterior
update still remains a Beta distribution. This allows for additive evidence integration, which is not only computationally more efficient,
but also can be extended to support an arbitrary set of data sources $S$. This aligns with the "Extensibility" goal of MERIT. The updates to the Beta
distribution are given by:
\begin{equation}
  \alpha_{total} = \alpha_{prior} + \sum_{s \in S} \alpha_s, \quad \beta_{total} = \beta_{prior} + \sum_{s \in S} \beta_s
\end{equation}

The final posterior distribution is:
\begin{equation}
  P(\theta \mid \text{Evidence}) \sim \text{Beta}(\alpha_{total}, \beta_{total})
\end{equation}

This probabilistic approach provides a robust framework for modelling uncertainty and resolving conflicts between
different evidence sources. For a verified signal, if a candidate has claimed evidence on a CV (low weight) and also
possesses the demonstrated evidence to back it up, such as on GitHub (high weight), then the resulting distribution will
skew towards a high probability of the candidate possessing the skill, with low variance. On the other hand, let's consider a skill 
that MERIT cannot verify. If that skill is listed heavily on their CV, but absent from their GitHub, the $\alpha$ parameter will only 
increase slightly. The baseline $\beta$ parameter will then prevent the score from artificially inflating, mathematically mitigating 
"keyword stuffing" even further.

Once we have computed the posterior distribution for a metric, we can then split it into its expected value and its variance. Effectively
this measures what the candidate's proficiency in a certain skill is, and the confidence that we have in that claim. For a random variable
$\theta \sim \text{Beta}(\alpha, \beta)$, its expected value $E[\theta]$ and variance $\text{Var}(\theta)$ are calculated as defined earlier in Equation~\ref{eq:beta_distribution}.

To summarise, the expected value of this distribution is the candidate's fused signal for the metric in question.
The scoring engine, as a result, grants higher scores to candidates who possess demonstrated evidence over candidates with only claimed evidence. 
However, this model is not strict enough to entirely rule out candidates who may have claimed evidence, but have no demonstrated evidence to back it up.
MERIT's Bayesian Fusion engine has a bypass for candidates who have no evidence signals for a given metric. Rather than falling back on the default $0.5$ expected
value from the prior, the engine intercepts the calculation, returning a flat score of $0.0$. This prevents candidates who do not provide evidence (claimed or demonstrated)
from ranking higher than candidates who do have evidence, albeit with scores below the prior's expected value.

\subsection{Glass-Box Payload}
The recent AI-driven innovations in ATS platforms have neglected explainability, turning towards
black-box approaches. Recruiters are often left without insight into how a candidate was ranked, they are unable to delve 
into the metrics that contribute to the candidate's score. MERIT's approach to address this is through a JSON payload that
can function as a mathematical audit trail. For any given metric, MERIT will show the full scoring process for it, from signal extraction,
to temporal decay, to Bayesian Fusion. When penalties, such as keyword stuffing, are applied, they are also raised in
this audit trail for recruiters to investigate further. 

This payload is generated during the execution stage and is structured to ensure the metrics that recruiters
observe can be broken down into its constituent parts. The payload schema consists of three main layers,
each layer gets more granular as we go down the hierarchy.

The first layer is the highest level layer, the aggregation layer, which contains the overall score and the
calculation summary. It provides a high-level overview and the global weighted formula used. 
The second layer is the middle layer, the metric layer. This contains the specific metrics
(e.g., \textit{Experience}, \textit{req\_Python}) stated in the job configuration. Each entry contains its independent
Bayesian parameters ($\alpha, \beta$) and its fused signal. The third layer is the evidence layer,
which is the lowest level layer, found within \texttt{source\_details}. This contains the
specific data source signals, confidence coefficients, and temporal decay penalties.

By exposing this structured data in a candidate's report, we justify the score for a given candidate: recruiters
can see how each data source contributed for a given metric, how Bayesian fusion combined those contributions into
a per-metric score, and how metrics are weighted by job priority to produce the final match score. Presenting every
stage in this form gives recruiters the means to challenge and interrogate the score rather than accepting it at
face value. Appendix~\ref{appendix:math_trail} shows a representative 
end-to-end audit trail in the UI for a given metric.
